\chapter{Arabic invoice AI pipeline}\label{ch:ai}

\section{Pipeline at a glance}

\begin{figure}[htbp]
\centering
\begin{tikzpicture}[
  node distance=6mm and 9mm,
  stage/.style={draw,rounded corners,minimum width=32mm,minimum height=11mm,align=center,font=\small,fill=orange!10},
  store/.style={draw,cylinder,shape border rotate=90,aspect=0.3,minimum width=22mm,minimum height=10mm,font=\small,fill=blue!7},
  arr/.style={-{Latex[length=2.2mm]},thick}
]
\node[store] (src) {PDF / e-mail};
\node[stage,right=of src]   (s1) {1.~Intake};
\node[stage,right=of s1]    (s2) {2.~OCR};
\node[stage,below=of s2]    (s3) {3.~Field\\extraction};
\node[stage,left=of s3]     (s4) {4.~Translation};
\node[stage,left=of s4]     (s5) {5.~Validation\\\& routing};
\node[store,below=of s5]    (out) {\texttt{structured/invoices/}};

\draw[arr] (src) -- (s1);
\draw[arr] (s1)  -- (s2);
\draw[arr] (s2)  -- (s3);
\draw[arr] (s3)  -- (s4);
\draw[arr] (s4)  -- (s5);
\draw[arr] (s5)  -- (out);
\end{tikzpicture}
\caption{The five pipeline stages; each emits a well-typed payload to
         the next. Only Stage~5 writes to the canonical corpus.}
\label{fig:pipeline}
\end{figure}

\section{Stage 1 -- Intake}

\paragraph{Inputs.} A source file, either:
\begin{itemize}
  \item a PDF dropped into a watched folder, or
  \item an e-mail attachment (e.g.\ the thread captured in
        \corpus\texttt{source/FW- AI~\_ Arabic Invoice\ldots.eml}).
\end{itemize}

\paragraph{Outputs.} A new entry staged for Stage~2 and a provisional
\texttt{source\_document} object:
\begin{lstlisting}[language=jsonx]
{ "path": "docs/arabic/source/<file>", "sha256": "<sha>" }
\end{lstlisting}
The SHA-256 is computed first and compared against
\corpus\texttt{structured/corpus.json}; a match short-circuits
re-processing (idempotence) and is the duplicate-payment control
(\S\ref{ch:controls}).

\paragraph{Precondition.} The MIME detector must classify the file as
one of \texttt{application/pdf}, the OOXML Word type
(\texttt{application/vnd.openxmlformats-\allowbreak officedocument.\allowbreak wordprocessingml.\allowbreak document}),
\texttt{text/plain}, or \texttt{message/rfc822}. All other types are
rejected with reason \textsc{unsupported\_media}.

\section{Stage 2 -- OCR}

\paragraph{Command.} The corpus has been produced with:
\begin{lstlisting}[language=bash,basicstyle=\ttfamily\small]
pdftoppm -r 300 <input.pdf> work/page -png
for img in work/page-*.png; do
  tesseract "$img" stdout -l ara+eng >> <output>.txt
done
\end{lstlisting}
The combined \texttt{ara+eng} language model is mandatory: invoices
routinely contain English brand names, Latin-script SKU codes, and
numerals alongside Arabic prose.

\paragraph{Quality gate.} Stage 2 emits a per-page confidence metric
(mean of tesseract's \texttt{conf} values over non-empty tokens). Pages
below \textbf{55\%} confidence are flagged for human review; below
\textbf{30\%} the file is rejected with reason \textsc{ocr\_low\_confidence}
and routed to manual processing.

\paragraph{Calibration.} The 30\,\%/55\,\% thresholds are not arbitrary.
They were calibrated empirically on the three-document Arabic golden
set: 30\,\% is the lowest per-page confidence observed on any page of
the ZATCA assessment that still yielded a schema-valid record after
human correction; 55\,\% is the lowest value at which Stage~3 achieved
zero field-extraction errors without human intervention. Both numbers
must be re-calibrated on the first ten invoices from each new country
added to \cref{tab:countries}.

\paragraph{Artefact.} The raw OCR text is written to
\corpus\texttt{extracted-text/<slug>.txt} and its path is recorded
in \texttt{source\_document.ocr\_text\_path}. This preserves the
training signal for future model refinement.

\section{Stage 3 -- Schema-guided field extraction}

\paragraph{Objective.} Produce a candidate record conforming to
\texttt{invoice.schema.json} from the OCR text. The extraction is
\emph{schema-guided}: the LLM prompt is constructed at run time by
serialising the schema (with \texttt{description} fields) so the
model is told exactly which JSON keys it may emit. This is the
single largest lever against hallucinated fields.

\paragraph{Prompt template.} Pseudocode; the concrete implementation
uses the Anthropic Messages API with structured-output tooling.

\begin{lstlisting}[language=bash,basicstyle=\ttfamily\small,caption={Schema-guided extraction prompt template.}]
SYSTEM: You extract tax-invoice fields into a JSON object that MUST
        validate against the provided JSON Schema. If a field is not
        present in the source, omit it -- NEVER invent a value.
USER:   <schema>
        ---
        OCR TEXT (Arabic may appear in RTL order; preserve spelling):
        <ocr_text>
        ---
        Emit JSON only. Set "metadata.language" to "ar", "en", or "ar+en"
        based on the dominant script.
\end{lstlisting}

\paragraph{Known-hard fields.} The following fields are the most error-prone
in the shipped corpus and must be handled with care:

\begin{table}[htbp]
\centering
\caption{Extraction pitfalls observed in the corpus and mitigations.}
\label{tab:pitfalls}
\small
\begin{tabularx}{\linewidth}{@{}l X X@{}}
\toprule
\textbf{Field} & \textbf{Pitfall} & \textbf{Mitigation} \\ \midrule
\texttt{electronic\_id}    & Tesseract maps the digit ``1'' to ``I'' in mixed alnum strings (e.g.\ \texttt{B1STE8}$\to$\texttt{BISTE8}). & Validate against the regex \verb|^[0-9A-Z]{26}$|; on mismatch, re-OCR at 600 dpi with \texttt{--psm~6}. \\
\texttt{totals.gross}      & Arabic digits (\ar{٠١٢٣٤٥٦٧٨٩}) may slip through unmapped. & Pre-tokenise: substitute Arabic-Indic digits with ASCII before parsing. \\
\texttt{vat\_rate\_percent}& Zero-rated vs exempt both show $0$. & Disambiguate with \texttt{vat\_category}: \textsc{zr} vs \textsc{ex} drawn from supplier's activity code lookup. \\
\texttt{signature.kind}    & ZATCA assessments have no signature block. & Default to \textsc{system\_generated} for \texttt{document\_kind} = \textsc{vat\_assessment}. \\
\texttt{seller.name\_en}   & Pure-Arabic invoices lack an English name. & Leave empty; Stage 4 will populate via translation. \\
\bottomrule
\end{tabularx}
\end{table}

\paragraph{Output.} A candidate record is written in memory only (never
persisted yet) and passed to Stage~4.

\section{Stage 4 -- Arabic$\to$English translation}

\paragraph{Scope.} This stage does \emph{not} translate the whole
document. It translates only the free-text fields that downstream AP
actually uses in English:

\begin{itemize}
  \item \texttt{seller.name\_en} and \texttt{buyer.name\_en} if absent;
  \item every \texttt{line\_items[].description\_en} if absent;
  \item \texttt{seller.address\_en} and \texttt{buyer.address\_en}.
\end{itemize}

Tax IDs, amounts, dates, and enum values are \emph{never} passed through
translation; they are copied verbatim.

\paragraph{Glossary.} A fixed glossary enforces terminology consistency.
Worked examples:

\begin{table}[htbp]
\centering
\caption{Locked glossary. The right-hand column is what the translation
         agent \emph{must} output.}
\label{tab:glossary}
\small
\begin{tabularx}{\linewidth}{@{}X X@{}}
\toprule
\textbf{Arabic} & \textbf{English (locked)} \\ \midrule
\ar{هيئة الزكاة والضريبة والجمارك} & Zakat, Tax and Customs Authority (ZATCA) \\
\ar{بنك ستاندرد تشارترد}           & Standard Chartered Bank \\
\ar{ضريبة القيمة المضافة}           & Value-Added Tax (VAT) \\
\ar{اشعار صدور فاتورة}              & Invoice Issuance Notification \\
\ar{فاتورة ضريبية}                  & Tax Invoice \\
\bottomrule
\end{tabularx}
\end{table}

\paragraph{Quality gate.} Back-translation is run as a sanity check:
the English output is translated back to Arabic and compared to the
source Arabic string. The metric is the \emph{normalised Levenshtein
distance} $d = \mathrm{lev}(a, b) / \max(|a|, |b|)$ computed at the
Unicode-codepoint level after Arabic NFKC normalisation (stripping
tatweel and harakat). The 0.30 ceiling was chosen because the three
golden-set Arabic strings have pairwise $d \leq 0.18$ between their
recorded Arabic and a human-verified back-translation, leaving a
$\geq 1.5\times$ safety margin. Distances above \textbf{0.30} are
flagged (\textsc{translation\_drift}) and the original Arabic is
retained as authoritative.

\section{Stage 5 -- Validation and routing}

\paragraph{Validation.} The candidate record is validated with
\texttt{validate.py} (\cref{ch:schema}). Any schema failure halts the
pipeline and the record is written to a side-queue
(\texttt{structured/invoices/rejected/}) with a \texttt{notes[]} entry
summarising the failure.

\paragraph{Routing.} On success, the pipeline consults the structured
\texttt{bh-accounts-payable-doi.json} and the \texttt{workflows/} set
to decide which state machine to instantiate:

\begin{enumerate}
  \item If the counterparty is a tax authority
        (\texttt{supplier\_type}~=~\textsc{government} or
         \texttt{tax\_authority}~=~\textsc{zatca})
        $\Rightarrow$ \texttt{vat-payment} workflow.
  \item Else if there is a PO reference in \texttt{notes[]}
        $\Rightarrow$ \texttt{po-eproc} workflow.
  \item Else if a prior \texttt{manual-recurring} record exists for
        the same \texttt{vendor\_id}
        $\Rightarrow$ \texttt{manual-recurring} workflow.
  \item Otherwise $\Rightarrow$ \texttt{manual-one-off} workflow.
\end{enumerate}

\paragraph{Commit.} The record is written to
\corpus\path{structured/invoices/<slug>.json} and a new item is
appended to \texttt{corpus.json}. This is the only stage in the
pipeline that mutates the corpus.

\section{Pipeline control plane}

\Cref{fig:pipeline-states} models the pipeline itself as a small state
machine; each invoice occupies exactly one of these states at any
moment. The control plane exposes three operations:
\texttt{ingest(file)}, \texttt{retry(id)}, \texttt{release(id)}.

\begin{figure}[htbp]
\centering
\begin{tikzpicture}[
  node distance=6mm and 5mm,
  s/.style={draw,rounded corners,minimum width=20mm,minimum height=7mm,align=center,font=\footnotesize,fill=gray!5},
  t/.style={-{Latex[length=2mm]}}
]
\node[s] (new)    {\textsc{received}};
\node[s,right=of new] (ocr) {\textsc{ocr\_ok}};
\node[s,right=of ocr] (ext) {\textsc{extracted}};
\node[s,right=of ext] (xlat) {\textsc{translated}};
\node[s,right=of xlat] (val)  {\textsc{validated}};
\node[s,below=of val] (route) {\textsc{routed}};
\node[s,below=of ocr] (hitl) {\textsc{hitl\_review}};

\draw[t] (new) -- (ocr);
\draw[t] (ocr) -- (ext);
\draw[t] (ext) -- (xlat);
\draw[t] (xlat) -- (val);
\draw[t] (val) -- (route);
\draw[t] (ocr) -- node[right,font=\tiny]{conf\textless 30\%} (hitl);
\draw[t] (val) to[bend left] node[above,font=\tiny,pos=0.3]{schema fail} (hitl);
\draw[t] (hitl) -- node[below,font=\tiny,pos=0.4]{corrected} (route);
\end{tikzpicture}
\caption{Pipeline control-plane states for a single invoice.}
\label{fig:pipeline-states}
\end{figure}

\section{Evaluation harness}

\paragraph{Golden set.} The three Arabic records in
\corpus\texttt{structured/invoices/} form the regression golden set.
Every code change must re-produce byte-identical records from the
source PDFs; \texttt{validate.py} plus a \texttt{jq}-driven equality
check is the verifier.

\paragraph{Per-field metrics.} For each round of Session-5 testing the
harness reports:

\begin{itemize}
  \item \textbf{Per-field exact-match accuracy} (over the 27 scalar
        fields in \texttt{invoice.schema.json}).
  \item \textbf{Per-field F1} on \texttt{line\_items[]} treated as sets
        (descriptions matched after glossary-aware fuzzing).
  \item \textbf{Human-approval rate.} Fraction of records released by
        HITL with zero edits. Target $\geq$~0.75 after Session~4,
        $\geq$~0.90 after Session~5.
\end{itemize}

\paragraph{Session acceptance.} The five-session model from
\cref{tab:sessions} is scheduled per country; the
\texttt{ai\_session\_phases} enum in \texttt{enums.yaml} is the system
of record. An AI rollout is complete for a country once the
Session-5 pass meets the human-approval-rate target against a cohort
of $\geq$~5 stakeholders.
